\documentclass[aps,pra,reprint,onecolumn,11pt,notitlepage,nofootinbib]{revtex4-2}
\usepackage{amsmath,amssymb,graphicx,booktabs,array}
\usepackage[colorlinks=true,linkcolor=blue,citecolor=blue,urlcolor=blue]{hyperref}
\newcommand{\norm}[1]{\left\|#1\right\|}
\newcommand{\eps}{\varepsilon}
\newcolumntype{L}[1]{>{\raggedright\arraybackslash}p{#1}}
\emergencystretch=2em
\begin{document}
\title{VigilODE polyalgorithm extensions:
Fused polynomial actions, finite-step error constants, and causal certification}
\author{VigilODE research supplement}
\noaffiliation
\date{28 October 2026}
\begin{abstract}
We investigate the six remaining development tasks following PR92, reusing
the existing conditioned resolvent, joint polynomial action, signed-adjoint
certificate, causal cache, and homotopy implementation. Three parallel theory
lanes produce seventeen concrete alternatives. The leading constructions are
a fused vector-coefficient Clenshaw action, a cubic midpoint defect bound with
constant $1/12$ and rational decay accumulation, variable-step two-channel
transport, and value-only continuation margins. We distinguish established
results, new project-specific derivations, finite exact-arithmetic checks,
and unimplemented floating-point obligations. Three programs, published before
execution, pass their declared exact-rational checks. On the stated targets,
fusion changes operator products from $5m$ to $m$, and the cubic planning
example changes its minimum dyadic grid from 512 to 256. These are operation
and planning results, not measured solver speedups. Runtime defaults, Rust
source, prior negative evidence, and unresolved deployment qualification are
preserved.
\end{abstract}
\maketitle

\section{Question, source identity, and useful reuse}
The inspected runtime baseline is PR92, commit
\par\noindent\texttt{7eb2dd0ffda39afb2871d76700a720986ff09fe0}.
\par The initial remote survey
found no later pull request. Task labels PR-N01 through PR-N06 denote planned
development tasks, not GitHub pull-request numbers. This supplement executes
the requested literature and constructive-exploration stages R2/R4 of
PhysMath Research Loop 1.3.0, with supporting derivations and checks. It does
not restart earlier ten-stage campaigns or their correction budgets.

Several proposed optimizations are already present. PR68 added joint actions
in Chebyshev and Laguerre bases; PR69 preserved an unsuccessful continuous-scaling
experiment. The latest code subsequently added signed output-adjoint
Laguerre certification and guarded Leja/Chebyshev/Laguerre routing.
The coupled solver already shares one projected factorization across eight
stages. Causal finite propagation and a nilpotent error majorant were already
proved in PR91. These are baselines to reuse, not new achievements here.

\begin{table}[h]
\caption{Existing capabilities and the proposed increment. Source paths are
relative to the repository's \texttt{crates/} directory.}
\begin{tabular}{L{0.29\linewidth}L{0.30\linewidth}L{0.30\linewidth}}
\toprule
Current source & Already available & Proposed increment\\
\midrule
\texttt{rodas5p-core/} polynomial action and router & Coefficient enclosures,
complete-total admission, exact cache keys, charged refusals & Fused-only
backward execution with its own local-error certificate\\
\texttt{rodas5p-integrators/} conditioned polynomial & Full-space relation,
conditioned chart, original residual & Prefer cheap structure; extend by a
proved perturbation bound\\
\texttt{vigilode-coupled/} causal and structural modules & Shared factor,
prefix reuse, represented-stage residual enclosure & Physical and
variable-step transport; carrier-specific endpoint budgets\\
\texttt{vigilode/} session and switcher & Scoped persistent session,
same-window work comparison & Retain winning session; conditional future
residence with an explicit cost-variation certificate\\
\bottomrule
\end{tabular}
\end{table}

The current guarded polynomial provider is an immutable \emph{dense symmetric}
operator. Its mathematical results do not automatically certify an opaque
matrix-free callback. Matrix-free, Jacobian-free, and derivative-free are
different properties. A model-declared generator can satisfy all three;
a supplied JVP remains derivative-assisted even without a stored matrix.

\section{Critical use of the supplied literature}
Robust numerical algebraic geometry studies parameter uncertainty and
exceptional algebraic strata\cite{robust}. Searching for nearby parameters
can clarify the structure of an uncertain polynomial system, but changing
an ODE parameter would change the original integration target. We reuse
explicit branch and parameter contracts, without importing parameter
regularization as an implicit acceptance rule.

The numerical nonlinear algebra survey distinguishes regular parameter
paths, root tracking, and certification\cite{nonlinear}. Its Newton-based
tracking and Jacobian-based Krawczyk method do not substantiate a
derivative-free implementation. A generic complex continuation theorem also
does not identify the prescribed real physical branch.

Kunkel--Mehrmann require regularity, derivative-array rank conditions,
coordinate choices and reconstruction for inherent ODEs\cite{dae}. Their
geometric symmetry construction is for linear time-varying DAEs. The paper
explicitly notes that it is inapplicable to its nonlinear pendulum example,
despite a self-adjoint Jacobian pair. A declared analytic reduction is a
useful future provider; generic RHS-only nonlinear DAE support remains open.

The ETH notes describe Hamiltonian composition and the Runge--Kutta
identities for invariants\cite{eth}. They motivate distinct structure-preserving
routes, while leaving adaptive selection and inexact-stage effects to be
analyzed separately.

\section{Fused backward polynomial actions}
For five distinct inputs and an output that requests only their sum, write
\begin{equation}
 F_m=\sum_{j=0}^{4}\sum_{k=0}^{m}c_{kj}P_k(X)w_j
     =\sum_{k=0}^{m}P_k(X)d_k,
 \qquad d_k=\sum_{j=0}^{4}c_{kj}w_j .
 \label{eq:fuse}
\end{equation}
The coefficient vectors can be formed on demand. The existing forward
implementation uses five vector recurrences for distinct inputs. Equation
\eqref{eq:fuse} permits one backward recurrence.

\subsection{Laguerre identity and error transfer}
Set $b_{m+1}=b_{m+2}=0$ and evaluate
\begin{equation}
 b_k=d_k+\frac{(2k+1)I-X}{k+1}b_{k+1}
           -\frac{k+1}{k+2}b_{k+2},\qquad k=m,\ldots,0.
 \label{eq:lag}
\end{equation}
Then $b_0=F_m$ for ordinary Laguerre polynomials. If each computed recurrence
has additive local defect $\delta_k$, including coefficient-vector formation,
operator-action error and rounded arithmetic, the exact identity is
\begin{equation}
 \widehat b_0-F_m=\sum_{k=0}^{m}L_k(X)\delta_k.
 \label{eq:lagdef}
\end{equation}
\textit{Proof.} The transfer of an impulse at position $k$ obeys the
Laguerre recurrence, with initial polynomials $1$ and $1-x$. It is therefore
$L_k(x)$. Superposition proves the action and perturbation identities,
including vector-valued coefficients. If $X$ is self-adjoint in the declared
Hilbert norm and $\sigma(X)\subset[0,L]$, the spectral theorem and DLMF
18.14.8\cite{dlmf} imply
\begin{equation}
 \norm{\widehat b_0-F_m}\le e^{L/2}\sum_k\norm{\delta_k}.
 \label{eq:lagbound}
\end{equation}
This bound is independent of degree in its transfer factor, but local
defects can grow with intermediate vectors. Large $L$ can make the bound
unhelpful. It is not valid to apply this factor to defects from the old
forward execution; the local defects depend on the evaluation direction.

For Chebyshev polynomials, use
\begin{equation}
 b_k=d_k+2Xb_{k+1}-b_{k+2}\quad(k\ge1),\qquad
 F_m=d_0+Xb_1-b_2.
\end{equation}
Telescoping gives the defect transfer $\sum_k T_k(X)\delta_k$. On a
self-adjoint represented chart contained in $[-1,1]$, its bound is
$\sum_k\norm{\delta_k}$. If rounding expands the interval, use the
corresponding $\cosh(k\operatorname{arcosh}(1+\eta))$ growth bound.

\subsection{Cost and admission}
At degree $m$, the fused-only distinct-input target changes operator products
from $5m$ to $m$. Coefficient-vector assembly still costs $O(5mn)$ scalar
operations. The current same-vector input already uses $m$ actions; callers
requesting all individual columns cannot use the shortcut. There is no
factor-five elapsed-time claim.

The new complete budget must include truncation, coefficient enclosures,
source combinations, recurrence defects, normalization, underflow and output
assembly. The current signed-adjoint Bernstein envelopes remain the
comparator. Their existing resource cap is not removed by an unimplemented
bound. A later matrix-free port needs immutable provider identity and a
certified action-error enclosure; any extra action or matrix storage used
for that certificate belongs in the cost.

\section{Resolvents, hybrids, and nonnormal alternatives}
\subsection{Exact Chebyshev degree}
For self-adjoint positive $A$, $\sigma(A)\subset[a,b]$, $0<a<b$, set
$c=(a+b)/2$, $d=(b-a)/2$, and
\begin{equation}
 p_m(z)=\frac{T_m((c-z)/d)}{T_m(c/d)},\qquad
 q_{m-1}(z)=\frac{1-p_m(z)}{z}.
\end{equation}
Since $p_m(0)=1$, $q_{m-1}$ is polynomial. Updating
$x_m=x_0+q_{m-1}(A)r_0$ gives $r_m=p_m(A)r_0$ and
\begin{align}
 \norm{r_m}&\le\delta_m\norm{r_0}, &
 \delta_m&=\frac{1}{T_m(c/d)}=\frac{2q^m}{1+q^{2m}},\\
 q&=\frac{\sqrt{b/a}-1}{\sqrt{b/a}+1}, &
 m_{\min}&=\left\lceil\frac{\operatorname{arcosh}(1/\eps)}
 {\operatorname{arcosh}(c/d)}\right\rceil .
\end{align}
The spectral theorem gives the norm bound. Alternation of $T_m$ gives the
minimax residual polynomial: a strictly better polynomial would have a
difference with $m$ roots between the extrema and an additional root at
zero outside the interval. This exceeds its degree.

An outward integer check must validate the chosen degree; rounded logarithms
alone are insufficient. For $[a,b]=[1,9]$, $q=1/2$ and the exact smallest
degree for $\delta_m\le10^{-8}$ is 28. If $a=b$, self-adjointness gives
$A=aI$, so a scalar solve is the appropriate route.

These methods eliminate orthogonalization and basis storage, but still
require dependent operator applications. Compare against CG, the existing
shared coupled basis, and specialized direct relations, not only a weak
GMRES implementation. Polynomial preconditioners can have useful fixed
linearity properties\cite{wathen}, while Laguerre short recurrences trade
inner products and storage against approximation behavior\cite{laguerre}.

\subsection{Reuse an approximate structural inverse}
Let $P$ satisfy $\norm{D}\le\theta<1$ for $D=I-AP$. Define
\begin{equation}
 Q_r=P\sum_{j=0}^{r-1}D^j .
 \qquad AQ_r=I-D^r,
 \qquad \norm{f-AQ_rf}\le\theta^r\norm{f}.
\end{equation}
The finite geometric sum proves this identity without requiring $AP=PA$.
For $A=A_0+E$, a sufficient bound is
$\theta\le\norm{I-A_0P}+\norm{E}\norm{P}$.
Existing exact exchange, block and quotient inverses can therefore serve as
cheap near-structure correction operators. Freeze the polynomial plan
within a standard linear outer solver, or use an explicitly flexible method.

\subsection{Nonnormality and numerical range}
For $N=\left(\begin{smallmatrix}0&1\\0&0\end{smallmatrix}\right)$,
$A=I+KN$, and $p(z)=1-z$, all eigenvalues equal one but
$\norm{p(A)}_2=|K|$. Eigenvalue placement alone does not justify the SPD
bound. Yet $A^{-1}=I-KN$ is exact: the same example favors the existing
declared low-degree relation route.

A verified numerical range permits the established constant
\cite{crouzeix}
\begin{equation}
 \norm{g(A)}_M\le(1+\sqrt2)\sup_{z\in W_M(A)}|g(z)|.
\end{equation}
If $C=(cI-A)/d$ has numerical range inside a filled Bernstein ellipse of
parameter $R$, with $1\le R<\zeta=\xi+\sqrt{\xi^2-1}$, $\xi=c/d>1$, then
\begin{equation}
 \norm{p_m(A)}_M\le(1+\sqrt2)
 \frac{R^m+R^{-m}}{\zeta^m+\zeta^{-m}}.
\end{equation}
The ellipse identity for $T_m$ and the maximum principle prove this bound.
It is a first Chebyshev/Faber branch requiring no general Faber engine.
General Leja/Faber choices need scalar and operator-domain error bounds
\cite{leja,faber}. Matrix-function experiments show that the best backend
depends on its integrator and problem\cite{dek}; a universal ordering is
not inferred.

\section{A finite-step cubic constant and a smaller justified grid}
Consider $z'=-\phi(z)$, $\phi(z)=kz+gz^3$, $k>0$, $g\ge0$, $z(0)=Y\ge0$.
Let $L(u)=k+3gu^2$, $D_h=1+hk/2$, and
$F_u(K)=K+h\phi(u+K/2)$. For the exact endpoints $u=z(t)$, $v=z(t+h)$,
\begin{equation}
 |F_u(v-u)|\le\frac{L(u)^2\phi(u)h^3}{12},\qquad
 |M_h(u)-v|\le\frac{L(u)^2\phi(u)h^3}{12D_h}.
 \label{eq:cubic}
\end{equation}
These are finite-step inequalities for every $h>0$.

\textit{Proof.} The exact trajectory is decreasing and convex, with
$0\le z''\le L(u)\phi(u)$. If $\ell(s)$ is its endpoint chord and
$m=(u+v)/2$, then
$0\le\ell(s)-z(t+s)\le L(u)\phi(u)s(h-s)/2$.
Define nonnegative quantities
\begin{align}
 A&=\int_0^h[\phi(\ell(s))-\phi(z(t+s))],ds,\\
 B&=\int_0^h\phi(\ell(s)),ds-h\phi(m)
   =\frac{ghm(u-v)^2}{4}.
\end{align}
The residual is $A-B$. Integrating the chord kernel gives
$A\le L(u)^2\phi(u)h^3/12$. Since $u-v\le h\phi(u)$ and
\begin{equation}
 L(u)^2-3gu\phi(u)=k^2+3kgu^2+6g^2u^4\ge0,
\end{equation}
the same upper bound holds for $B$. Use $|A-B|\le\max(A,B)$.
Finally $F_u'(K)\ge D_h$ on the real line proves the endpoint bound.
Derivatives enter this analytic proof, not the runtime callback contract.

The midpoint map is globally nonexpansive because
\begin{equation}
 M_h'(x)=\frac{1-(h/2)\phi'(m)}{1+(h/2)\phi'(m)},\qquad |M_h'(x)|\le1.
\end{equation}
Thus numerical error is transported separately from the exact trajectory
used in the local truncation term. Since
$z_j\le Ye^{-kjh}\le Yq^j$, $q=(1+kh)^{-1}$, and the polynomial
$L(x)^2\phi(x)$ has nonnegative coefficients, a cumulative bound is
\begin{equation}
 \mathcal T_N=\frac{h^3}{12D_h}\left[
 k^3Y S_N(q)+7k^2gY^3S_N(q^3)
 +15kg^2Y^5S_N(q^5)+9g^3Y^7S_N(q^7)\right],
 \label{eq:sum}
\end{equation}
where $S_N(r)=\sum_{j<N}r^j$. For dyadic $N$, the updates
$(p,s)\mapsto(p^2,s(1+p))$, starting with $(r,1)$, compute this sum with
$O(\log N)$ positive operations. No transcendental library is needed.

\begin{table}[h]
\caption{Registered exact-rational planning example:
$k=1$, $g=1/2$, $Y=T=1$, $\eps=10^{-5}$. Displayed decimals are diagnostic;
admission comparisons were exact.}
\begin{tabular}{lrr}
\toprule
Bound & Truncation bound at $N=256$ & Minimum dyadic $N$\\
\midrule
Inherited constant $1/8$ & $1.7846536915\,10^{-5}$ & 512\\
Constant $1/12$ only & $1.1897691277\,10^{-5}$ & 512\\
$1/12$ and rational decay & $3.3853071299\,10^{-6}$ & 256\\
\bottomrule
\end{tabular}
\end{table}

The remaining budget must cover root error, update rounding and clock
displacement through the complete prefix gate. An admissible residual
threshold alone does not pay all these terms. The inherited Picard fallback
also retains its contraction gate. A separate bracketed monotone backend
could relax that backend-specific restriction, but must not silently call
an unqualified noncontractive fallback.

\section{Structural and causal error transport}
For $P=11^T/n$, $Q=I-P$, $J=-dQ-(d-\kappa)P$, $d\ge\kappa\ge0$, a discrete
step has $S_k=\sigma_{k,t}Q+\sigma_{k,c}P$. Fixed projectors permit
\begin{equation}
 \prod_{\ell>k}S_\ell=p_{k,t}Q+p_{k,c}P,
 \qquad p_{k,r}=\prod_{\ell>k}\sigma_{\ell,r}.
\end{equation}
At each $h_k$, the existing scalar stage matrices $L_k(r)$ give adjoints
$w_{k,r}=L_k(r)^{-T}\beta$. A full-stage residual $\rho_{k,i}$ contributes
$(w_{k,t,i}Q+w_{k,c,i}P)\rho_{k,i}$ to the endpoint. The exact induced
infinity norm of a two-channel block is
\begin{equation}
 \nu_n(t,c)=\left|t+\frac{c-t}{n}\right|
       +(n-1)\left|\frac{c-t}{n}\right|.
\end{equation}
Unrolling the error recurrence and applying this norm gives the
variable-step residual transport, with sharp independent-block row sums.
Temporal defects, endpoint arithmetic and clock discrepancy remain separate
terms. Products over an accepted grid are an a posteriori statement, not a
certificate for an unknown future adaptive grid.

For the declared affine model $y'=b+Jy$, the exact flow is
\begin{align}
 F_\Delta(y)={}&e^{-d\Delta}Qy+e^{-(d-\kappa)\Delta}Py\\
 &+\Delta\varphi_1(-d\Delta)Qb
  +\Delta\varphi_1(-(d-\kappa)\Delta)Pb.
\end{align}
If $\zeta_k=\widehat y_{k+1}-F_{\Delta_k}(\widehat y_k)$ is enclosed,
\begin{equation}
 E_{k+1}\le e^{-(d-\kappa)\Delta_k}E_k+\norm{\zeta_k}_\infty.
\end{equation}
The exponential has nonnegative entries and row sum
$e^{-(d-\kappa)t}$, proving the scalar factor. This also suggests a direct
flow route when it is cheaper than executing and certifying eight stages.
Validated exponential/$\varphi$ arithmetic is still required.

For the general raw-$U$ causal equations,
\begin{equation}
 WU_i=h\gamma f\left(y+\sum_{j<i}a_{ij}U_j\right)
       +\gamma\sum_{j<i}c_{ij}U_j,\qquad W=I-h\gamma J_0,
\end{equation}
let $\norm{W^{-1}}\le\omega$ and let $L_i$ bound the RHS Lipschitz constant
on the relevant stage domains. Then
\begin{equation}
 e\le\omega\sum_{\ell=0}^{s-1}B^\ell\rho,
 \qquad B_{ij}=\omega(h\gamma L_i|a_{ij}|+|\gamma c_{ij}|),\quad j<i.
\end{equation}
Subtract the stage equations to obtain $e\le\omega\rho+Be$; strict lower
triangularity gives $B^s=0$. This instantiates the already-proved causal
majorant. It does not need a global contraction factor below one.

For a carrier-aware affine stage matrix with inverse block bounds $D$ and
nonlinear remainder majorant $N$, the sharper form is
$e\le\sum_{\ell<s}(DN)^\ell D\rho$.
Endpoint weights convert this to a state error budget. A changed carrier's
temporal or stiff-uniform order remains a separate question.

\section{Continuation, session reuse and temporal residence}
For $G_\lambda=(1-\lambda)a+\lambda G$ on a closed ball
$\overline B(c,r)$, suppose $G$ is $q$-Lipschitz and certified quantities
satisfy $e\ge\norm{G_\lambda(c)-c}$, $b\ge\norm{G(c)-a}$. The next
station $\mu=\lambda+\delta\le1$ leaves the ball invariant if
\begin{equation}
 (\lambda+\delta)qr+e+\delta b\le r.
\end{equation}
This follows by adding and subtracting $G_\mu(c)$ and $G_\lambda(c)$.
When the denominator is nonzero and the margin positive,
\begin{equation}
 \delta\le\frac{(1-\lambda q)r-e}{qr+b},\qquad \mu q<1
\end{equation}
provides a value-only Banach station and residual-to-root constant
$1/(1-\mu q)$. The strict contraction check is separate. If $q\ge1$, this
does not establish the endpoint $\lambda=1$. If the endpoint is already
contractive, intermediate stations may add unnecessary cost.

For a monotone scalar residual, a certified bracket and
$F(v)-F(u)\ge m(v-u)$ give uniqueness and error $|F(x)|/m$ without
runtime derivatives. Small residual alone cannot choose between the two
roots of $x^2-p=0$.

Session reuse has a different role. Preserving every state component read
by a deterministic integrator, its callback bindings, output landings and
arithmetic order implies pause/resume trace equality by induction. This
requires FSAL validity, controller history, buffers, cumulative budgets and
epochs; a $(t,y,h)$ snapshot is insufficient. The first candidate is to
continue a winning probe inside the existing scoped session rather than
introducing a self-referential owned callback design.

Finally, same-window cost evidence does not establish future residence.
Under an explicit future variation certificate
$C_B(0)-C_A(0)\ge d>0$ and
$|(C_B-C_A)(j)-(C_B-C_A)(j-1)|\le V$, telescoping gives
\begin{equation}
 C_B(j)-C_A(j)\ge d-jV,\qquad
 \sum_{j=1}^{m}(C_B(j)-C_A(j))\ge md-\frac{Vm(m+1)}2.
\end{equation}
Require $mV<d$ and the latter bound greater than switching/calibration cost
$S$. For continuous rates with Lipschitz constant $L>0$, the corresponding
bound is $dH-LH^2/2$ for $H\le d/L$. If $L=0$, all $H\ge0$ preserve the
advantage and strict payback is $dH>S$.
Smooth ODE coefficients do not supply such a work certificate: adaptive
acceptance thresholds can cause discrete jumps. Trace-margin certification
and its cost remain a separate research task.

\section{Geometric alternatives and two exact counterexamples}
Declared separable Hamiltonians admit force-only splitting without
Jacobian/JVP/Newton evaluations. For oscillator Verlet,
\begin{equation}
 M(h)=\begin{pmatrix}1-z^2/2&h\\
 -h\omega^2(1-z^2/4)&1-z^2/2\end{pmatrix},\qquad z=h\omega.
\end{equation}
The determinant is one, and the fixed-step power-bounded interval is
$0<h\omega<2$. The endpoint has a nontrivial Jordan block. Yet
$\omega=1$, $h_1=1$, $h_2=19/10$ obey the individual interval and satisfy
\begin{equation}
 \det(M(h_2)M(h_1))=1,\qquad
 \operatorname{tr}(M(h_2)M(h_1))=-9661/4000<-2.
\end{equation}
The periodic product is unstable. Moreover, the exact free-particle flow
with state-dependent $h(q,p)=\epsilon q$ has determinant $1+\epsilon p$,
so state-dependent selection can destroy symplecticity itself.

An inexact RK stage can nevertheless have a useful invariant-drift
certificate. Let $I(y)=y^TCy$, $C=C^T$, and $y^TCf(y)=0$.
For $F_i=f(Y_i)+\eta_i$, stage defects
$\delta_i=Y_i-y-h\sum_j a_{ij}F_j$, and $z=y+h\sum_i b_iF_i$,
direct quadratic expansion gives
\begin{align}
 I(z)-I(y)={}&-h^2\sum_{ij}(b_i a_{ij}+b_j a_{ji}-b_i b_j)F_i^TCF_j\nonumber\\
 &-2h\sum_i b_i\delta_i^TCF_i+2h\sum_i b_iY_i^TC\eta_i.
\end{align}
For a symplectic RK tableau the first term vanishes. Endpoint error $r$
adds $2z^TCr+r^TCr$. This bounds invariant drift; it does not make an
inexact input-output map exactly symplectic.

Other concrete branches in the catalog are an exact small invariant block
plus a polynomial complement, predetermined affine prefix scans, and a
declared regular DAE reduction provider. They retain their own structural
and reconstruction hypotheses.

\section{Verification, development priorities and remaining obligations}
The three checker programs and their protocol were published before execution:
\par\noindent\texttt{866668348f840bd43558859d2795576e292b6ca7}.
\par
All passed. Polynomial checks cover degrees 0--8 and all impulse positions,
the exact degree-28 example and a Jordan counterexample. Homotopy checks
include 108 cubic algebra fixtures, exact grid decisions, continuation-ball
and branch controls, and conditional residence. Four geometry checks compare
scalar adjoints with independently assembled full block matrices, verify a
nonlinear causal majorant with norm greater than one, check invariant drift,
and reproduce the Verlet counterexamples.

These checks do not replace the written proofs. They do not validate new
binary64 kernels or measure native performance. Earlier exploratory CAS
checks are explicitly identified; the later registration does not apply
retroactively to them. One bounded cross-lane review corrected the explicit
$L=0$ wording in the continuous residence statement. No checker source or
runtime arithmetic changed after the registered run.

\begin{table}[h]
\caption{Six task contracts. All native implementations remain pending.}
\begin{tabular}{L{0.10\linewidth}L{0.38\linewidth}L{0.39\linewidth}}
\toprule
Task & First implementation candidate & Remaining theoretical/validation bridge\\
\midrule
N01 & Fused-only backward Clenshaw; retain current coefficient/admission code &
New local-defect enclosure, complete cost, matrix-free authority; nonnormal witness\\
N02 & $1/12$ rational-decay outward planner & Prefix/update/clock budget,
minimum-grid and range-boundary oracles\\
N03 & Variable-step two-channel transport; direct-flow comparator & Physical
defect enclosure and charged adaptive consumer\\
N04 & Retain winning scoped session and workspace capacity & Full trace
identity, generic domain and branch certificate\\
N05 & Raw-$U$ residual-to-endpoint budget and work/span comparison & Nonlinear
closure, carrier order, domain and stiff-uniform accuracy\\
N06 & Distinguish observed cost from certified variation & Counterfactual
future trace/variation certificate and its acquisition cost\\
\bottomrule
\end{tabular}
\end{table}

The refreshed machine-readable plan preserves 23 inherited nodes, 34 edges,
eight native suites, seven unresolved qualification tasks and existing
correction budgets. Seventeen candidate records refine the six tasks rather
than creating seventeen new mandatory PRs. The deployment sequence remains
release qualification, matched-accuracy hardware measurements, and actual
client qualification. Earlier Givens, fair-AB, SUNDIALS, lint/API, negative
control and authority-reference issues are not declared resolved here.

The polyalgorithm first filters candidates by target, structure, derivative
permission and physical certificate. Only then does it compare total setup,
actions, memory, verification, retries and failed work. Hardware service
times are measured quantities; exact numerical constants cannot determine
a universal fastest policy. Theory can fix the admissible region and error
budget while local evidence decides the economic choice within that region.

Complete proofs, candidate contracts, exact outputs, source-bound receipts,
the Korean report and local Codex handoff are stored in
\texttt{research/polyalgorithm\_alternatives\_20261028/}. The unchanged
PR92 evidence retains its original scope. This supplement supplies a concrete
implementation basis without claiming a general nonlinear closure theorem
or a new production speedup.

\begin{thebibliography}{99}
\bibitem{robust} E. R. Cobian, J. D. Hauenstein, and C. W. Wampler,
\textit{Robust Numerical Algebraic Geometry},
\href{https://arxiv.org/pdf/2403.18749}{arXiv:2403.18749}.
\bibitem{nonlinear} D. J. Bates et al., \textit{Numerical Nonlinear Algebra},
\href{https://arxiv.org/html/2302.08585v2}{arXiv:2302.08585v2}.
\bibitem{dae} P. Kunkel and V. Mehrmann, \textit{Discretization of inherent
ODEs and the geometric integration of DAEs with symmetries},
\href{https://arxiv.org/pdf/2205.08105}{arXiv:2205.08105}.
\bibitem{eth} H. Ammari, \textit{Geometrical numerical integration methods
for differential equations}, ETH Lecture 5,
\href{https://people.math.ethz.ch/~grsam/FS18/NAII/files/slides/ODE-Lecture5.pdf}{lecture notes}.
\bibitem{dlmf} NIST Digital Library of Mathematical Functions,
\href{https://dlmf.nist.gov/18.14.E8}{Eq.18.14.8} and
\href{https://dlmf.nist.gov/18.12.E13}{Eq.18.12.13}.
\bibitem{wathen} A. J. Wathen and T. Rees,
\textit{Chebyshev semi-iteration in Preconditioning}, Oxford Report08/14,
\href{https://www.cs.ox.ac.uk/files/1540/NA-08-14.pdf}{author report} (2008).
\bibitem{laguerre} B. N. Sheehan, Y. Saad, and R. B. Sidje,
\textit{Computing exp(-tau A)b with Laguerre polynomials}, ETNA
\textbf{37}, 147--165 (2010).
\bibitem{crouzeix} M. Crouzeix and C. Palencia,
\textit{The numerical range as a spectral set},
\href{https://arxiv.org/pdf/1702.00668}{arXiv:1702.00668} (2017).
\bibitem{leja} M. Caliari, P. Kandolf, A. Ostermann, and S. Rainer,
\textit{The Leja method revisited: backward error analysis for the matrix exponential},
\href{https://arxiv.org/pdf/1506.08665}{arXiv:1506.08665} (2016).
\bibitem{faber} B. Beckermann and L. Reichel,
\textit{Error estimation and evaluation of matrix functions via the Faber transform},
\href{https://pro.univ-lille.fr/fileadmin/user_upload/pages_pros/bernhard_beckermann/abs/paper15.pdf}{author manuscript}.
\bibitem{dek} P. J. Deka, M. Tokman, and L. Einkemmer,
\textit{A comparison of Leja- and Krylov-based iterative schemes for Exponential Integrators},
\href{https://arxiv.org/pdf/2211.08948}{arXiv:2211.08948} (2022).
\end{thebibliography}
\end{document}
